% Compile with LuaLaTeX
\documentclass[11pt,a4paper]{article}

% ==========================================
% 1. MODERN FONTS & LANGUAGES
% ==========================================
\usepackage{fontspec}
% Tip: For Romance languages and IPA, "Charis SIL" or "Doulos SIL" are the best fonts. 
% If you have them installed, uncomment the next line:
% \setmainfont{Charis SIL} 

\usepackage{polyglossia}
\setmainlanguage{english} % change to german if needed
% Load the Romance languages you study
\setotherlanguages{french,spanish,italian,portuguese,romanian,latin}

% ==========================================
% 1b. TYPOGRAPHY & QUOTES
% ==========================================
\usepackage{microtype}   % Better line breaks, especially helpful for multilingual docs
\usepackage{csquotes}    % Language-aware quotes: \enquote{...} picks the right marks
\usepackage{booktabs}
\usepackage{graphicx}

% ==========================================
% 2. LINGUISTICS & PHONETICS
% ==========================================
\usepackage{tipauni}     % Modern, LuaLaTeX-safe IPA symbols
\usepackage{gb4e}        % For interlinear glosses (standard in linguistics)
\noautomath              % CRITICAL: Prevents gb4e from crashing with modern packages
\usepackage{forest}      % For drawing modern syntax trees

% ==========================================
% 3. MATH & LAYOUT
% ==========================================
\usepackage{amsmath,amssymb,mathtools}
\usepackage{geometry}
\geometry{margin=2.5cm}
\usepackage{hyperref}
\usepackage{xcolor}
\usepackage{siunitx}



\title{Computational Linguistics \\ \large Romance Languages Template}
\author{Your Name}
\date{\today}

\begin{document}
\maketitle

\section{Multilingual Text}
Because you work with Romance languages, you can type native characters
directly. \texttt{polyglossia} handles the hyphenation rules automatically,
and \texttt{csquotes} picks the correct quotation marks for whichever
language is currently active --- just write \verb|\enquote{...}|.

\begin{itemize}
    \item \textbf{English:} \enquote{The student ate an apple.}
    \item \textbf{French:} \textfrench{L'élève a mangé une pomme. \enquote{Bonjour.}} (é, è, ê, ç, œ --- and French guillemets with the correct narrow space inside)
    \item \textbf{Spanish:} \textspanish{El niño comió una manzana. \enquote{¡Hola!}} (ñ, á)
    \item \textbf{Portuguese:} \textportuguese{A criança comeu uma maçã. \enquote{Olá.}} (ã, ç)
    \item \textbf{Romanian:} \textromanian{Copilul a mâncat un măr. \enquote{Bună.}} (ă, â, î)
\end{itemize}

\section{Phonetics and IPA}
You can use \texttt{tipauni} shortcuts or just type Unicode symbols directly. 
\begin{itemize}
    \item French \textit{pain}: \textipa{[pɛ̃]}
    \item Portuguese \textit{são}: \textipa{[sɐ̃w̃]}
    \item Spanish \textit{perro}: \textipa{[pero]} vs \textit{pero}: \textipa{[peɾo]}
\end{itemize}

\section{Interlinear Glosses (\texttt{gb4e})}
This is the core tool for morphological analysis. Here is how a standard Romance gloss looks:

\begin{exe}
  \ex \gll L' élève a mangé une pomme. \\
            the student has eaten an apple \\
      \glt `The student ate an apple.' 
      \label{ex:french}
  \ex \gll El niño comió una manzana. \\
            the boy ate an apple \\
      \glt `The boy ate an apple.' \hfill (Spanish)
\end{exe}

\section{Syntax Trees (\texttt{forest})}
For syntax classes, \texttt{forest} is much more powerful and modern than the old \texttt{qtree} or \texttt{tikz-qtree}. Here is a simple phrase structure tree:

\begin{center}
\begin{forest}
  [CP
    [C' [C $\emptyset$]
      [IP
        [DP [D el] [N niño]]
        [I' [I comió]
          [VP [V \_\_\_] [DP [D una] [N manzana]]]
        ]
      ]
    ]
  ]
\end{forest}
\end{center}

\section{Computational Tool: Dependency Parsing}
Since you are in \textit{computational} linguistics, you might need to represent dependency grammar (common in NLP tools like spaCy or UDPipe). While there are specific packages for this (like \texttt{tikz-dependency}), you can also use standard tables for basic alignment:

\begin{center}
\begin{tabular}{|c|c|c|}
\hline
\textbf{Token} & \textbf{Lemma} & \textbf{Dependency Relation} \\
\hline
El & el & det \\
niño & niño & nsubj \\
comió & comer & root \\
una & uno & det \\
manzana & manzana & obj \\
\hline
\end{tabular}
\end{center}

\section{Pro-Tip for ComputLing}
If you need to paste code, data structures (like JSON), or Python code into your linguistic papers, use the \texttt{listings} package (or \texttt{minted} for Pygments-powered highlighting). Here is a Python snippet analysing Romance text:

\begin{verbatim}
import spacy
# Load Spanish NLP model
nlp = spacy.load("es_core_news_sm")
doc = nlp("El niño comió una manzana.")
for token in doc:
    print(token.text, token.pos_, token.dep_)
\end{verbatim}

\end{document}